\documentclass[11pt]{article}
\usepackage[a4paper,margin=2.5cm]{geometry}
\usepackage{fontspec}
\setmainfont{Noto Serif}
\usepackage{setspace}\onehalfspacing
\usepackage{graphicx}
\usepackage[hidelinks]{hyperref}
\usepackage{enumitem}
\usepackage[super,sort&compress]{natbib}
\usepackage{xcolor}
\setlength{\parskip}{4pt}
\newcommand{\authors}[1]{\textcolor{red}{[#1]}}
\title{Supplementary Information}\date{}
\begin{document}
\maketitle
\emph{Stratospheric warmings shift winter cold risk rather than creating two kinds of event.} Supplementary Notes hold text moved from the main text to keep it within the article length; every number is from \texttt{ssw-design-analysis/\allowbreak{}CONSOLIDATED\_RESULTS.md} or the result JSONs it indexes. Supplementary Tables are built by \texttt{10\_tables/\allowbreak{}tableS*.py} from \texttt{results/\allowbreak{}current/\allowbreak{}}.

\section*{Supplementary Note 1 | Event differences: the bound, its assumption and its falsifier}

How much the forced response differs between events can be estimated only under assumptions. In CMIP6 the forced variance between events, after removing the variance already present before onset, is 0.019σ² [−0.033, 0.072]: the forced probability of a negative outcome ranges across events from about 0.62 to 0.87 (0.47 to 0.93 at the upper bound), accounting for 2\% of the variance of a single event's label (at most 8\%). Any one label is therefore almost entirely the shift plus chance. The bound assumes that an event's forced response is uncorrelated with the internal variability it adds to. It also gives a falsifier: at its upper limit, forced differences could produce at most about half the class contrast of the criterion in CMIP6 and two thirds in observations, so, under the latent-effect model tested, classes producing those full contrasts would have to exceed the measured, assumption-dependent bound; subtypes with smaller effects are not excluded. An audit finds the falsifier robust to the strongest damping of internal variability the experiment allows and to forcing that depends on the pre-onset state; it fails only if forcing anti-correlated with concurrent internal variability (ρ $\le$ −0.26 to −0.40), which no available data can test (Methods).

\section*{Supplementary Note 2 | Archetypes, post-onset coupling and vortex geometry}

The field's archetype pair shows the same (Fig. 1c): February 2018 and January 2019 received comparable forced shifts and similar forced probabilities of a downward outcome (0.80 and 0.75 across nine models), and at the primary initialisations the observed outcomes lie within the central 95\% of the model distributions in 17 of 18 model–event cases (Extended Data Fig. 2), although they are conventionally taken as opposite kinds. The lower stratosphere after onset does track how much events differ (week-2 100 hPa anomaly against the later surface response, r = 0.85 across 18 ensembles of distinct events\cite{loeffel26}; Methods), but between members of one SNAPSI ensemble, where the forcing is shared, the relation is weak and equally present with no SSW (r = 0.11 against 0.09; Extended Data Fig. 3), and in CMIP6 the post-onset stratosphere predicts the surface as well on ordinary winter days: it is a general property of stratosphere–troposphere coupling, not a feature that sorts SSWs into kinds. The difference most often proposed, vortex geometry, cannot be tested in SNAPSI or our zonal-mean CMIP6 fields, so we test it in reanalysis\cite{seviour13}. Events classified as vortex splits tend towards a stronger early surface response than displacements\cite{lehtonen16,nebel24}, but not significantly: −0.25σ [−0.61, +0.13] over days 8–25 (p = 0.28; 37 events; Methods). With 27 splits and 10 displacements the interval still admits a split response stronger by 0.6σ, as expected when a 1,000-year simulation needs more than 50 events to separate the two\cite{maycock15}.

\section*{Supplementary Note 3 | Regional residuals beyond the polar-cap circulation}

The polar-cap index is not the whole of the regional response: high-latitude Europe and mid-latitude East Asia cool more than the circulation implies, and mid-latitude North America warms (+0.49σ) where the circulation relation implies cooling, with the same sign of residual in observations (+0.55 K, p = 0.09), in line with the mixed North American signal during weak-vortex states\cite{kretschmer18,huang21} (Methods). What an SSW changes regionally is therefore a distribution of the regional variable itself.

\section*{Supplementary Note 4 | Operational discrimination}

Operational forecasts agree on event differences: in seven subseasonal systems, pre-onset skill for the 500 hPa polar-cap response after 13 SSWs lay at the 91st percentile of that at random winter dates, not significantly higher\cite{nebel24}. In ten systems of the S2S reforecast archive\cite{vitart17} (ECMWF as a discovery set, nine others confirmatory; Methods), forecasts started 2–9 days before onset track event-to-event differences in the surface response no better than on event-free dates of the same season and lead (ECMWF r = 0.77 against 0.52, p = 0.11, 12 SSWs; the nine others 0.23 against 0.52, p = 0.91, 17 SSWs; Fig. 5a,b), which rules out large event-specific skill, not small (Methods).

\section*{Supplementary Note 5 | Why the operational bias did not recur}

Exploratory checks added afterwards point to the events, not the systems: the newer ECMWF cycle ranks the earlier SSWs much as the older one did (paired difference +0.019 [−0.055, +0.076]); given the kind of outcome, the periods agree (upward outcomes ranked 0.76 in 1998–2021 and 0.75 in 2023–24); what changed is the mix, 13 of 17 earlier outcomes downward against one of four since; and the signs of the 2023–24 outcomes are 3.8 times likelier under the forecasts as issued than under the 1998–2021 bias (Methods). Whether models under-represent the surface response to SSWs is disputed\cite{sigmond13,kolstad20,afargan24,dai25}; our results bound such a bias as one of the size of the shift, of uncertain persistence, rather than of event-to-event differences.

\section*{Supplementary Note 6 | Continuity at the wind reversal with ERA5 winds, 1940–2025}

The NCEP test of the main text uses 114 weakening episodes of 1958–2025. Its local regression-discontinuity estimates were unstable (four of 16 placebo cutoffs with p < 0.05). Several NCEP episodes just above zero are reversals in other reanalyses, so measurement error sits exactly at the cutoff. The ERA5 version (registered before the ERA5 data were retrieved) uses the ERA5 10 hPa, 60° N wind from 1940. Its minima agree with NCEP's (r = 0.98 on 96 shared episodes), but five lie on the other side of zero: December 1958 (NCEP −5.8, ERA5 +6.4 m s⁻¹), January 1968 (+1.1, −4.7), January 1977 (0.0, −2.6), March 1981 (+0.9, −0.7) and February 2002 (+1.8, −0.4). The dose response replicates for all four outcomes. No outcome shows a step after correction. The largest step, −0.73 on the Arctic Oscillation (raw p = 0.03; Holm 0.11), is the strongest hint of a step at reversal anywhere in this work. In CMIP6, 5,726 episodes bound the step at +0.05σ [−0.02, +0.13]. In ERA5, reversals follow larger decelerations (balance p = 0.002), so a step estimated at zero partly reflects how much faster the wind fell. The local estimates again fail their placebo cutoffs, so the observational conclusion rests on the global estimates, as registered. All numbers are in Methods and in \texttt{results/\allowbreak{}current/\allowbreak{}5\_mechanism/\allowbreak{}vortex\_threshold\_continuity.json} (key \texttt{era5\_1940\_2025}).

\section*{Supplementary Note 7 | The shift rule out of sample}

The 39 events of 1959–2022 verify the SNAPSI probabilities, but only the predictor was out of sample. Twelve SSWs of 1941–1958 in ERA5 provided a test whose outcomes had not been seen. The rule failed it at the registered 10th-percentile threshold (0 of 12 against 0.32; p = 0.012).

The early events differ in three ways:

\begin{itemize}[leftmargin=*,itemsep=1pt]
  \item \textbf{They are weaker reversals.} The median minimum wind is −3.8 m s⁻¹, against −8.5 m s⁻¹ in 1959–2022.
  \item \textbf{The data are older.} ERA5's stratosphere in the 1940s rests on few soundings reaching 10 hPa.
  \item \textbf{There are only twelve of them.}
\end{itemize}

None of these explains the failure on its own. The dose response accounts for about a quarter of the smaller cooling, and the failure persists from 1946, when the upper-air record improves.

Two readings follow, and with twelve events they cannot be separated:

\begin{itemize}[leftmargin=*,itemsep=1pt]
  \item the single SNAPSI probability, which rests on two strong imposed events, overstates the risk after weak reversals;
  \item the 1941–1958 sample was unusually mild.
\end{itemize}

Either way, a per-SSW probability verified on one set of events did not transfer to another. The forecast product should therefore scale with the size of the disturbance and be verified prospectively. Full counts are in Supplementary Table 4.

\section*{Supplementary Note 8 | Members that do and do not reverse the vortex from one initial state}

The ten S2S reforecast systems give a comparison with SNAPSI's design, in which every member shares one initial state and members either have an SSW or do not. Here the SSW is not imposed: some members reverse the 10 hPa wind on their own. The test was registered before any split was computed.

Of 461 starts with at least two members of each kind, the onset comes early enough for the 8–25-day window in 40 (162 for the 8–20-day window). Results are in Methods.

The signature of a threshold on one population appears:

\begin{itemize}[leftmargin=*,itemsep=1pt]
  \item \textbf{Rate:} members that reverse are more often downward.
  \item \textbf{Contrast:} their class contrast is unchanged.
  \item \textbf{Spread:} they are no more spread out.
  \item \textbf{Shape:} their lower tail moves more than their upper. The quantile shift is −0.58σ at the 5th percentile against +0.12σ at the 95th, the asymmetry also seen in CMIP6.
\end{itemize}

Two results are weaker than in SNAPSI:

\begin{itemize}[leftmargin=*,itemsep=1pt]
  \item \textbf{The mean shift is small} (−0.15σ over the shorter window).
  \item \textbf{Within starts there is no dose relation and no step.} This matches SNAPSI's own finding that, between members sharing one forcing, the post-onset stratosphere predicts the surface only weakly.
\end{itemize}

Spontaneous reversals in these ensembles mostly come late in the forecast (only 40 of the 461 mixed starts have the median onset by lead 9). The comparison therefore adds evidence on the form of the change (translation, unchanged contrast), not on its size.

\section*{Supplementary Note 9 | Extended results: the label as a threshold (rates, contrasts, criterion sweep)}

Applying the surface conditions of the Karpechko criterion to every member splits each ensemble into downward (DW) and non-downward (NDW) members (adapted to the length of the runs; in observations we also apply it as published; Methods). What the SSW changes is the rate: 84\% of members are DW with the SSW imposed against 45\% without it (Fig. 1a; 79\% against 41\% in the Southern Hemisphere). The same holds in observations: the share of the 39 observed events with ERA5 outcomes that meet the surface conditions, 69\% (54\% with all three published conditions), is reproduced by event-free dates displaced by the measured shift (75\% [61, 87]; Fig. 1b), so “about two thirds” is what one shifted population produces. In the Northern Hemisphere the contrast between the classes does not change: −1.62σ with the SSW imposed and −1.59σ without it (paired difference −0.03σ [−0.12, +0.05] on 25 centre–initialisation pairs; Fig. 2a), while the rate separates the arms in every model (Fig. 2b). Without an SSW the contrast is what a cut at zero gives on one Gaussian population; with it, members without an SSW that are shifted by the imposed effect and classified identically reproduce it (Methods). In ERA5, event-free dates reproduce 97\% of the surface DW–NDW contrast of the published criterion (93\% of that of an 850 hPa variant\cite{lu26}). Nor does this depend on how the criterion is set: across 108 versions varying its window, level, thresholds and stratospheric condition, the observed downward rates are those of shifted event-free dates (calibrated joint test p = 0.85, and 0.75 with the stratospheric condition; two planted classes 1.8σ apart are detected with probability 0.82), and those dates reproduce 76–154\% of the class contrast (Extended Data Fig. 4).

What has not been tested directly is whether the label, or the warming itself, identifies a kind of event. Three questions are easily conflated: whether downward and non-downward outcomes are distinct kinds of response or one distribution cut by a threshold; how much events differ in the response they force; and whether knowing the event improves a forecast. Events can differ without forming classes. The downward label differs from the split–displacement typology in a way that matters for testing it: it is defined on the surface outcome itself, so its classes differ by construction, and a test needs an intervention that fixes the forcing and a null that does not depend on how the criterion is set. We also ask what the shift means for regional cold risk and what forecasts get wrong. Attribution studies in the same experiment quantified the stratosphere's role in individual extremes\cite{seviour26}; our subject is the classification. We use an experiment that imposes the same observed SSW on every ensemble member in nine models (SNAPSI\cite{hitchcock22}), 1,517 SSWs in 20 CMIP6 simulations of 10 models, 43 observed SSWs and 114 observed vortex weakenings, and reforecasts and real-time forecasts of ten operational systems. Extended Data Table 1 lists every test and whether it was fixed before its data.

\section*{Supplementary Note 10 | Extended results: continuity at the wind reversal}

The same reasoning applies one step earlier, without a model. An SSW is defined by a threshold, the reversal of the zonal wind at 10 hPa, 60° N. If reversal marked a distinct kind of event, surface outcomes would jump at zero among otherwise similar vortex disruptions. In a test registered before any outcome was related to the wind, across 114 vortex weakenings since 1958, 42 of them SSWs, the outcome instead grows smoothly with how far the wind falls (Extended Data Fig. 6): per 10 m s⁻¹ weaker minimum wind the Arctic Oscillation over days 8–52 falls by 0.36 [0.21, 0.51], northern Eurasia cools by 0.36 K [0.07, 0.66] and a downward label becomes 0.08 [0.01, 0.14] more likely, with no step at reversal for any outcome (Arctic Oscillation −0.28 [−0.99, +0.47]). With 114 episodes this excludes only steps about as large as the whole SSW effect, and local discontinuity estimates fail their placebo checks (Methods). In CMIP6 the same rule yields 5,726 episodes and a step of +0.05σ [−0.02, +0.13], whose upper 90\% limit is 15\% of the 0.8σ shift observed after SSWs, against a dose of 0.17σ per 10 m s⁻¹. With ERA5 winds for 1940–2025 instead (144 weakenings, 54 of them reversals), which place five NCEP episodes on the other side of zero, the dose replicates (Arctic Oscillation 0.29 [0.15, 0.43] per 10 m s⁻¹; northern Eurasia 0.38 K) and no outcome shows a step after correction for the four tested; the largest, on the Arctic Oscillation, is −0.73 [−1.51, −0.06] (Holm-adjusted p = 0.11), and reversals there also follow larger decelerations (Supplementary Note 6). The tests resolve no surface-response discontinuity at reversal within these limits; “SSW” behaves here as a threshold on a continuum, which does not deny that reversal may matter dynamically.

\section*{Supplementary Note 11 | Extended results: the shift rule as a forecast}

The shift provides a fixed, model-derived probabilistic benchmark; on these 39 events it verifies. In a test registered before it was run, the probability of a cold northern-Eurasian period after an SSW was taken from the nudged experiment alone (0.24, 0.32 and 0.45 for periods below the event-free 5th, 10th and 20th percentiles) and verified, untuned, on the 39 observed SSWs: the observed frequencies, 0.21, 0.26 and 0.38, are consistent with those probabilities (p = 0.40–0.71) and not with climatology (p = 0.0006–0.008), and the risk ratio rises with severity, as a shift implies (observed 1.9, 2.6 and 4.1 from the 20th to the 5th percentile; Extended Data Fig. 7). A negative annular mode, given 0.74 by the CMIP6 events, followed 30 of 43 observed SSWs (climatology 0.42). A user who protects whenever the risk exceeds their cost/loss ratio recovers 60–75\% of the value of a perfect forecast at ratios near the climatological rate, and loses value between the observed rate and the slightly higher model probability (Methods). Only the predictor is out of sample: the observed frequencies had been reported before. The verification holds without the two SSWs that SNAPSI imposes (37 events: p = 0.38–0.57 under the rule, 0.002–0.012 under climatology) and at every threshold from the 33rd to the 2.5th percentile, and it is bracketed by a second rule built independently from the CMIP6 circulation shift and the ERA5 relation between circulation and temperature on event-free dates (0.11, 0.19 and 0.33), which under-predicts where SNAPSI slightly over-predicts (Supplementary Tables 1 and 3). The product has limits that it must carry: over high-latitude Europe the models overstate the most extreme cold (none of 39 events below the event-free 2.5th percentile, against 20\% predicted); over mid-latitude East Asia the observed rise is not distinguishable from climatology with 39 events, and over North America neither models nor observations show raised cold risk; and in weeks 3–4, where SNAPSI rests mainly on the later initialisations of its two events, northern-Eurasian cold was rarer than predicted (0.15 against 0.27) and not distinguishable from climatology (Supplementary Table 2). Out of sample the rule did not verify: after twelve SSWs of 1941–1958 in ERA5, whose surface outcomes no one had examined, no northern-Eurasian period fell below the event-free 10th percentile, against 32\% predicted (p = 0.012), and cold risk did not rise detectably above that period's climatology (Supplementary Table 4). These reversals were weaker (an exploratory check; median minimum wind −3.8 against −8.5 m s⁻¹ in 1959–2022), but by the dose response that explains only about a quarter of their smaller mean cooling (0.2 against 1.0 K), and 1940s stratospheric winds are weakly constrained; a single probability per SSW therefore holds for the events it was verified on and over-predicted for this earlier, weaker set.

\section*{Supplementary Note 12 | Extended results: event information and operational skill}

If the event information available at onset determined how strongly an event couples, adding it should improve a forecast of that coupling; limited gains with these predictors and methods do not bound physical predictability. We compare two probabilistic forecasts of each CMIP6 event's surface response (days 8–52): the shifted distribution alone (SSW responses in other models, by calendar date) and an event-aware forecast that adds the stratospheric state up to onset. Scored with the continuous ranked probability score on events of held-out models, the event-aware forecast improves on the shift by 0.8\% [−0.3, 1.4] (Fig. 4a). The same information improves forecasts on size-matched ordinary winter days by 3.7\% [2.3, 5.3], and in none of 200 such samples was it worth as little as after SSWs; the two-component model gains 1.1\% [0.4, 2.0], against 3.1\% on ordinary days (a comparison added afterwards). The result holds in five robustness checks (Methods). The realised label cannot be an issue-time predictor, because it is defined from the later outcome (a model can still predict its probability); the test is of the event information available at onset. Point prediction agrees: the event adds −0.006 [−0.033, +0.020] to the out-of-sample R² of the surface response before onset (Fig. 4b) and +0.015 [−0.028, +0.059] with the post-onset stratosphere, when 96\% of that diagnostic skill is present without an SSW (Fig. 4c; Methods).

Operational forecasts agree on event differences\cite{nebel24}: in ten systems of the S2S reforecast archive\cite{vitart17}, forecasts started 2–9 days before onset track event-to-event differences in the surface response no better than on event-free dates of the same season and lead (Fig. 5a,b; Supplementary Note 4), which rules out large event-specific skill, not small.

\section*{Supplementary Note 13 | Extended results: operational forecasts in 1998–2021 and after}

Over 1998–2021, what the forecasts got wrong was the size of the shift. In ECMWF, 4\% of observed outcomes fell in the two outer rank bins against 17\% expected (a non-symmetric rank histogram does not by itself exclude omitted structure). They fall on the negative-NAM side, as registered before the confirmatory test was run: mean rank 0.40 against 0.54 on event-free dates (p = 0.005; Holm-adjusted over all primary tests, 0.03), lower than the event-free value in nine of ten systems (Fig. 5c), a deficit specific to SSWs that survives the further checks in Methods. The systems gave a negative-NAM period after these SSWs a mean probability of 0.57, whereas it occurred 82\% of the time, and ensemble-mean polar-cap responses were about half those observed (116 against 241 Pa). Observed northern-Eurasian temperature lay on the cold side in all ten systems, not significantly after correction (mean rank 0.39 against 0.53, p = 0.014; Holm-adjusted 0.07; Fig. 5d). The error was not merely a missed warming: starts whose ensembles reversed the 10 hPa wind near the observed onset under-predicted as much as those that did not (difference +0.003 [−0.05, +0.05]). A follow-up registered before the 100 hPa forecasts were retrieved places the loss in the lower stratosphere: the forecast 100 hPa anomaly was too weak (p = 0.005) while the 10 hPa wind was not, and given the forecast 100 hPa state the surface was not detectably under-predicted (p = 0.30; Methods) — the lower-stratospheric persistence deficit reported for nearly all systems\cite{garfinkel25}, not a coupling deficit.

The bias did not recur. After the three SSWs of 2023 and 2024, held out and registered before they were scored, and after the SSW of 4 March 2026, scored with real-time forecasts of five systems under a design registered before any 2026 forecast was retrieved, each observed polar-cap outcome lay, on average over systems, on the upward side of the ensembles (four events: mean rank 0.70 against 0.55, p = 0.93; Methods, Extended Data Fig. 5). Exploratory checks added afterwards point to the events, not the systems: given the kind of outcome, the periods agree, and what changed is the mix, 13 of 17 earlier outcomes downward against one of four since (Supplementary Note 5). Whether models under-represent the surface response to SSWs is disputed\cite{sigmond13,kolstad20,afargan24,dai25}; our results bound such a bias as one of the size of the shift, of uncertain persistence, rather than of event-to-event differences.

\section*{Supplementary Note 14 | Is the imposed SSW a translation? Residual quantiles against a declared tolerance}

A pooled variance ratio near 1 and a non-rejected shape test do not show that a distribution is translated, so the translation was tested directly. In every nudged/control pair the nudged quantiles were compared with the control quantiles moved by the difference in means, against a tolerance fixed before the test: ±0.25σ in any quantile and ±0.05 in a tail probability.

Results:

\begin{itemize}[leftmargin=*,itemsep=1pt]
  \item \textbf{Polar-cap annular mode, all 36 ensembles pooled:} every interval falls inside the tolerance. This is a positive equivalence result, not a failure to reject.
  \item \textbf{Single events:} unresolved. In February 2018 the lower tail is compressed: the most negative states are less extreme than a pure translation implies.
  \item \textbf{Northern-Eurasian temperature:} unresolved. There is a small excess of extreme cold below the 5th percentile, inside the tolerance, and lower-quantile intervals that cross it.
\end{itemize}

So translation is established for the circulation index at the resolution of these ensembles. For regional temperature it is neither established nor refuted. Figures are in Methods; values per pair are in \texttt{results/\allowbreak{}current/\allowbreak{}8\_experiment/\allowbreak{}snapsi\_residual\_quantiles.json}.

\section*{Supplementary Note 15 | Does the shifted null transport? Whole-winter cross-validation}

The shift used to build the null had been estimated on the same events it was compared with. Here it is estimated without each held-out winter; climatology and the constant class rate are also estimated without that winter. The test is retrospective, because the 39 outcomes had been seen, but nothing about a held-out winter enters its own prediction.

Results:

\begin{itemize}[leftmargin=*,itemsep=1pt]
  \item \textbf{Calibrated label rates out of fold:} 27 downward events observed against 29.2 expected; no version of the criterion is miscalibrated.
  \item \textbf{Better than climatology for the outcome itself} (CRPS).
  \item \textbf{No better from a constant class rate:} the label carries no information beyond the shift, as one shifted population predicts.
\end{itemize}

Values are in Methods and in \texttt{results/\allowbreak{}current/\allowbreak{}9\_literature/\allowbreak{}criterion\_transport\_cv.json}.

\section*{Supplementary Note 16 | How much do event-conditioned responses differ? The 18 ICON ensembles}

Loeffel et al. (2026) re-ran 18 SSWs from a long ICON simulation as 40-member spin-off ensembles. Their public archive gives the ensemble mean and spread of the standardised polar-cap geopotential anomaly. It does not give member trajectories, so only the differences between the ensemble-mean responses can be assessed. Distribution shape and label rates cannot.

The test was registered before any value was read. Onsets are where the ensemble-mean 10 hPa wind first turns negative. The variance of the ensemble means between events, net of finite-ensemble noise, is bounded by assuming the window-mean spread lies between zero and the daily spread.

Results:

\begin{itemize}[leftmargin=*,itemsep=1pt]
  \item \textbf{Days 8–25:} between-event variance 0.41–0.43, about 80\% of the within-ensemble variance.
  \item \textbf{Days 26–42:} 0.10–0.12, about 13–15\%.
  \item \textbf{Link to the lower stratosphere:} the event means track the week-2 100 hPa anomaly (r = 0.80).
\end{itemize}

These conditional means differ greatly from the CMIP6 forced-variance estimate (0.019) because they estimate different things:

\begin{itemize}[leftmargin=*,itemsep=1pt]
  \item \textbf{They keep each event's initial tropospheric state,} which persists into the first weeks. The fall from days 8–25 to days 26–42 is consistent with that. The CMIP6 estimate subtracts the differences present before onset.
  \item \textbf{The 18 events were selected,} so their spread is not a population variance.
\end{itemize}

What survives is the distinction the main text now draws: realised responses to different SSWs differ substantially and in proportion to the lower-stratospheric anomaly. How much of that the stratosphere itself forces is not settled by these summaries or by the assumption-dependent CMIP6 bound.

Values are in \texttt{results/\allowbreak{}current/\allowbreak{}5\_mechanism/\allowbreak{}icon\_event\_heterogeneity.json}.

\section*{Supplementary Note 17 | Forecasting with the continuous state, and against operational ensembles}

The paper's fixed SSW rules give one probability per event. The alternative tested here, registered before it was run, is a forecast that uses only what is known on the issue day and is told nothing about SSWs. It is a Gaussian whose mean depends on:

\begin{itemize}[leftmargin=*,itemsep=1pt]
  \item the 10 hPa wind, its 10-day change and its 10-day minimum;
  \item the surface annular mode and regional temperature over the previous five days;
  \item the season.
\end{itemize}

It was trained on every winter day of the training winters. Everything was fitted inside whole-winter folds and evaluated on all 57 SSWs that ERA5 records since 1940. Because those outcomes had been examined before, this is retrospective cross-validation.

Results:

\begin{itemize}[leftmargin=*,itemsep=1pt]
  \item \textbf{The state model scores best of the statistical forecasts.} It improves on climatology and on the fixed SNAPSI rule by about 10\%, in both eras and with any winter left out. The intervals include zero, so by the registered rule this is not a forecasting result.
  \item \textbf{The SSW label adds nothing.} Adding an indicator that the state is an SSW does not help. Neither does knowing how far the wind would fall over the next 20 days (the oracle). For this regional target, what an SSW adds is carried by the continuous state already present on the issue day.
  \item \textbf{The fixed SNAPSI rule over-predicts cold over the whole record}: 18.1 cold periods expected against 13 observed.
  \item \textbf{Dynamical ensembles started within a week after onset beat every statistical forecast decisively.} Routine bias and spread calibration improves them further, by 8\% (5–11\%).
\end{itemize}

The forecasting conclusion is therefore modest: after an observed SSW, the useful product is the calibrated dynamical ensemble. The statistical results add one conceptual point: being an SSW carries no extra information once the continuous state is known.

Values per event are in \texttt{results/\allowbreak{}current/\allowbreak{}6\_predictability/\allowbreak{}state\_forecast.json}.

\section*{Supplementary Note 18 | Further limits}

The main text keeps its limits short; these qualifications apply as well.

\begin{itemize}[leftmargin=*,itemsep=1pt]
  \item The operational tests rest on 4 to 20 events per system or set, and errors are shared within a winter, so the effective sample is closer to the number of winters.
  \item Regional samples of 39 events exclude only residuals larger than about a quarter of a standard deviation.
  \item The two-component model's statistical components need not correspond to physical mechanisms, and overlapping mechanisms are not excluded by a unimodal fit.
  \item The test of whether the realised lower-stratospheric dose accounts for the state-dependent shape had no useful power (Methods, Response shape and dose).
  \item Translation is established for the polar-cap annular mode pooled over ensembles; for single events and regional temperature the residual-quantile intervals are wider than the tolerance (Supplementary Note 14).
  \item The forecast comparisons and the transport test of the shifted null are retrospective cross-validations of outcomes that had been examined before; only the twelve SSWs of 1941–1958 and the March 2026 event were scored without prior inspection.
  \item The CMIP6 bound on forced differences assumes that an event's forced response is uncorrelated with the internal variability it adds to; it fails only if forcing anti-correlates with concurrent internal variability (Supplementary Note 1).
  \item Local regression-discontinuity estimates at the wind reversal fail their placebo checks in both reanalyses, so the continuity result rests on the global estimates, which exclude only steps about as large as the whole SSW effect.
\end{itemize}

{\small\bibliographystyle{unsrt}\bibliography{refs}}
\end{document}
